%! AP Statistics — Unit 1, Lesson 1.6 — Slides (Beamer)
\documentclass[aspectratio=169, 11pt]{beamer}
\usepackage{apstats-beamer}
\usepackage[most]{tcolorbox}

%=============================================================================
\begin{document}

% ─────────────────────────────────────────────────────────────────────────────
% SLIDE 1 — LESSON TITLE
% ─────────────────────────────────────────────────────────────────────────────
{
\setbeamercolor{background canvas}{bg=navy}
\begin{frame}[plain]
  \begin{tikzpicture}[remember picture, overlay]
    \fill[navylight] (current page.north west)
      rectangle ([yshift=-1.35cm]current page.north east);
  \end{tikzpicture}
  \vspace{2.8cm}
  \hspace{0.55cm}%
  \begin{minipage}{0.9\textwidth}
    {\color{goldacc}\bfseries\small LESSON 1.6}\\[0.5cm]
    {\color{white}\bfseries\LARGE Describing the Distribution of a\\[0.25cm]
      Quantitative Variable}\\[1.4cm]
    {\color{skymid}\small AP Statistics~$\cdot$~Unit 1: Exploring One-Variable Data}
  \end{minipage}
\end{frame}
}

% ─────────────────────────────────────────────────────────────────────────────
% SLIDE 2 — WARM-UP
% ─────────────────────────────────────────────────────────────────────────────
{
\setbeamercolor{background canvas}{bg=sky}
\begin{frame}[t, plain]
  \navyheader{Warm-Up}
  \vspace{1em}%
  \hspace{0.55cm}%
  \begin{minipage}{0.9\textwidth}
    \small

    A histogram shows homework hours per week for 30 students:\\[6pt]
    \quad 0--4 hrs: 3 students \quad 4--8 hrs: 9 \quad
    8--12 hrs: 11 \quad 12--16 hrs: 6 \quad 16--20 hrs: 1

    \vspace{0.6em}
    \begin{enumerate}
      \setlength{\itemsep}{0.7em}\setlength{\topsep}{0em}
      \item Describe the \textit{shape} of this distribution. Use one of
            these words: symmetric, skewed right, skewed left, bimodal,
            uniform. Explain your choice.
      \item Estimate a \textit{typical} (center) homework time for these
            30 students. What supports your estimate?
      \item Are there any \textit{outliers} --- values far from the overall
            pattern? Which interval(s)?
    \end{enumerate}

    \vspace{0.5cm}
    {\color{navy}\itshape\small
      Today we formalize these three observations into a four-part AP
      framework called \textbf{SOCS}: Shape, Outliers, Center, Spread.}
  \end{minipage}
\end{frame}
}

% ─────────────────────────────────────────────────────────────────────────────
% SLIDE 3 — PRIMARY OBJECTIVE
% ─────────────────────────────────────────────────────────────────────────────
{
\setbeamercolor{background canvas}{bg=sky}
\begin{frame}[t, plain]
  \begin{tikzpicture}[remember picture, overlay]
    \fill[navy] (current page.north west)
      rectangle ([yshift=-0.25cm]current page.north east);
  \end{tikzpicture}
  \vspace{0.45cm}\par
  \hspace{0.55cm}%
  \begin{minipage}{0.9\textwidth}

    \sectionlabel{Primary Objective}

    \normalsize
    Students will describe the distribution of a quantitative variable
    using the \textbf{SOCS} framework: \textbf{S}hape, \textbf{O}utliers,
    \textbf{C}enter, and \textbf{S}pread. Students will recognize common
    shapes and relate shape to real-world contexts.

    \vspace{0.45cm}

    \normalsize\textbf{\color{navy}By the end of this lesson, I will be able to:}

    \smallskip
    \begin{itemize}
      \setlength{\itemsep}{0.3em}\setlength{\topsep}{0.2em}
      \item Identify and name the shape of a distribution from a graph.
      \item Write a complete SOCS description with context in every sentence.
      \item Explain why outliers matter and how they affect center and spread.
      \item Predict the relationship between mean and median based on shape.
    \end{itemize}

  \end{minipage}
\end{frame}
}

% ─────────────────────────────────────────────────────────────────────────────
% SLIDE 4 — HOOK: TWO HISTOGRAMS
% ─────────────────────────────────────────────────────────────────────────────
{
\setbeamercolor{background canvas}{bg=hookbg}
\begin{frame}[t, plain]
  \begin{tikzpicture}[remember picture, overlay]
    \fill[goldacc] (current page.north west)
      rectangle ([yshift=-1.35cm]current page.north east);
    \node[anchor=west, text=white, font=\bfseries\large,
          inner xsep=0.55cm, inner ysep=0pt]
      at ([yshift=-0.68cm]current page.north west)
      {Hook --- Two Histograms, Two Very Different Shapes};
  \end{tikzpicture}
  \vspace{1.2cm}\par
  \vspace{0.6em}%
  \hspace{0.55cm}%
  \begin{minipage}{0.9\textwidth}
    \small

    % ── Two rough histogram sketches drawn with TikZ ─────────────────────────
    \begin{tikzpicture}[x=0.9cm, y=0.9cm]

      % ── Histogram A (left) — income, skewed right ──────────────────────────
      \node[white, font=\small\bfseries, anchor=south] at (2.5, 3.3)
            {Histogram A: U.S.\ Household Income};
      \draw[white, thick] (0, 0) -- (5.2, 0);
      \draw[white, thick] (0, 0) -- (0, 3.3);
      % bars: very tall at left, shrinking to right
      \fill[goldacc!80] (0.1,0) rectangle (0.9,3.0);
      \fill[goldacc!80] (1.0,0) rectangle (1.8,1.8);
      \fill[goldacc!80] (1.9,0) rectangle (2.7,0.9);
      \fill[goldacc!80] (2.8,0) rectangle (3.6,0.4);
      \fill[goldacc!80] (3.7,0) rectangle (4.5,0.15);
      % x-label
      \node[white, below, font=\scriptsize] at (2.5,-0.1) {Income};
      % shape label
      \node[white, font=\bfseries\small] at (2.8, 2.5) {Skewed right?};

      % ── Histogram B (right) — heights, symmetric ───────────────────────────
      \begin{scope}[xshift=6.8cm]
        \node[white, font=\small\bfseries, anchor=south] at (2.5, 3.3)
              {Histogram B: Heights of Adult Women};
        \draw[white, thick] (0, 0) -- (5.2, 0);
        \draw[white, thick] (0, 0) -- (0, 3.3);
        % bars: short, tall, tallest, tall, short (bell)
        \fill[skymid!80] (0.1,0) rectangle (0.9,0.6);
        \fill[skymid!80] (1.0,0) rectangle (1.8,1.8);
        \fill[skymid!80] (1.9,0) rectangle (2.7,3.0);
        \fill[skymid!80] (2.8,0) rectangle (3.6,1.8);
        \fill[skymid!80] (3.7,0) rectangle (4.5,0.6);
        \node[white, below, font=\scriptsize] at (2.5,-0.1) {Height (in)};
        \node[white, font=\bfseries\small] at (2.8, 2.5) {Symmetric?};
      \end{scope}

    \end{tikzpicture}

    \vspace{0.35cm}
    {\color{white}\small
      \textbf{Discuss:} Describe each histogram in one word. What real-world
      reason explains why income has that shape? Why do heights look different?}
  \end{minipage}
\end{frame}
}

% ─────────────────────────────────────────────────────────────────────────────
% SLIDE 5 — VOCABULARY: SHAPE
% ─────────────────────────────────────────────────────────────────────────────
{
\setbeamercolor{background canvas}{bg=greenbg}
\begin{frame}[t, plain]
  \navyheader{Vocabulary --- Shape}
  \vspace{0.8em}%
  \hspace{0.55cm}%
  \begin{minipage}{0.9\textwidth}
    \small
    \begin{description}[leftmargin=0em, itemsep=0.5em, font=\bfseries\color{navy}]
      \item[Symmetric] Roughly mirror-image about the center; mean $\approx$ median.
            \textit{Example: heights, IQ scores.}
      \item[Skewed right (positive skew)] Long tail toward \textit{higher} values;
            bulk of data at the left. Mean $>$ median.
            \textit{Example: household incomes, house prices.}
      \item[Skewed left (negative skew)] Long tail toward \textit{lower} values;
            bulk of data at the right. Mean $<$ median.
            \textit{Example: age at retirement, scores on an easy test.}
      \item[Unimodal] One clear peak.
      \item[Bimodal] Two clear peaks --- often signals two subgroups in the data.
      \item[Uniform] All bars roughly equal height --- no clear peak.
    \end{description}

    \vspace{0.3em}
    {\color{navy}\bfseries\small Key:}
    {\small The direction of skew = direction the \textbf{tail} points,
    \textit{not} the peak.}
  \end{minipage}
\end{frame}
}

% ─────────────────────────────────────────────────────────────────────────────
% SLIDE 6 — THE SOCS FRAMEWORK
% ─────────────────────────────────────────────────────────────────────────────
{
\setbeamercolor{background canvas}{bg=white}
\begin{frame}[t, plain]
  \navyheader{The SOCS Framework}
  \vspace{0.8em}%
  \hspace{0.55cm}%
  \begin{minipage}{0.9\textwidth}
    \small

    \begin{tabularx}{\textwidth}{@{} >{\bfseries\color{navy}\large}c
                                      >{\bfseries\color{navy}}p{0.15\linewidth}
                                      X @{}}
      S & Shape    & Overall pattern: symmetric, skewed right/left, unimodal,
                     bimodal, uniform. \\[0.5em]
      O & Outliers & Values far from the overall pattern. Identify and
                     investigate; do not automatically remove. \\[0.5em]
      C & Center   & Informal estimate of a typical value (median or balance
                     point). Exact measures come in Lesson 1.7. \\[0.5em]
      S & Spread   & Variability; for now, describe the range (min to max).
                     Formal measures in Lesson 1.7. \\[0.5em]
    \end{tabularx}

    \vspace{0.45cm}
    \begin{tcolorbox}[colback=sky, colframe=navy, boxrule=0.7pt, arc=1.5mm,
        left=2mm, right=2mm, top=1.5mm, bottom=1.5mm]
    \small\textbf{AP Context Rule:} Every sentence in a SOCS description must
    name the \textbf{variable} and the \textbf{population}. A sentence without
    context cannot earn full AP credit.
    \end{tcolorbox}

  \end{minipage}
\end{frame}
}

% ─────────────────────────────────────────────────────────────────────────────
% SLIDE 7 — MODEL SOCS RESPONSE
% ─────────────────────────────────────────────────────────────────────────────
{
\setbeamercolor{background canvas}{bg=white}
\begin{frame}[t, plain]
  \navyheader{Model SOCS Response --- SAT Math Scores ($n = 200$)}
  \vspace{0.8em}%
  \hspace{0.55cm}%
  \begin{minipage}{0.9\textwidth}
    \small

    \textit{Histogram description: symmetric, unimodal; center near 520;
    range 300--750; two students above 750.}

    \vspace{0.6em}

    \begin{tcolorbox}[colback=sky, colframe=navy, boxrule=0.7pt, arc=1.5mm,
        left=3mm, right=3mm, top=2mm, bottom=2mm]
    \small
    \textbf{[S]} The distribution of SAT Math scores for these 200 students
    is approximately \textit{symmetric and unimodal}, with no obvious skew.\\[4pt]
    \textbf{[O]} There appear to be two \textit{potential outliers} above 750
    points, which stand out from the rest of the distribution.\\[4pt]
    \textbf{[C]} The \textit{center} of the distribution is approximately
    520 points, representing a typical SAT Math score for these students.\\[4pt]
    \textbf{[S]} The \textit{spread} is roughly 300 to 750 points, a range
    of about 450 points.
    \end{tcolorbox}

    \vspace{0.4em}
    {\color{navy}\small\bfseries What earns full credit:}
    \begin{itemize}
      \setlength{\itemsep}{0.2em}
      \item All four SOCS components present.
      \item Context (\textit{SAT Math scores for these 200 students}) in every sentence.
      \item Outliers \textit{hedged} (``appear to be'') --- no formal rule yet.
      \item Units included with center (points) and spread (points).
    \end{itemize}
  \end{minipage}
\end{frame}
}

% ─────────────────────────────────────────────────────────────────────────────
% SLIDE 8 — SKEW AND MEAN VS. MEDIAN
% ─────────────────────────────────────────────────────────────────────────────
{
\setbeamercolor{background canvas}{bg=white}
\begin{frame}[t, plain]
  \navyheader{Shape Affects Which Summary to Use}
  \vspace{0.8em}%
  \hspace{0.55cm}%
  \begin{minipage}{0.9\textwidth}
    \small

    \begin{tabularx}{\textwidth}{@{} >{\bfseries\color{navy}}p{0.22\linewidth}
                                       X
                                       >{\bfseries\color{navy}}p{0.22\linewidth} @{}}
    \rowcolor{sky}
    \textbf{Skewed Right} & & \textbf{Skewed Left} \\
    \midrule
    Tail $\to$ right & &  Tail $\to$ left \\
    Mean $>$ Median  & &  Mean $<$ Median \\
    \multicolumn{3}{c}{\textit{The mean is pulled in the direction of the tail.}} \\
    \midrule
    \rowcolor{sky} \multicolumn{3}{c}{\textbf{Symmetric: Mean $\approx$ Median}} \\
    \end{tabularx}

    \vspace{0.55cm}
    {\color{navy}\bfseries\small Why does this matter?}

    \vspace{0.3em}
    \begin{itemize}
      \setlength{\itemsep}{0.4em}
      \item For skewed data, the \textbf{median} is a better center because it
            is resistant to outliers and extreme values.
      \item The \textbf{mean} is appropriate for symmetric distributions with
            no outliers.
      \item On the AP exam, knowing the relationship between shape and center
            measure is a frequently tested concept (Lesson 1.7 makes this precise).
    \end{itemize}
  \end{minipage}
\end{frame}
}

% ─────────────────────────────────────────────────────────────────────────────
% SLIDE 9 — GUIDED PRACTICE
% ─────────────────────────────────────────────────────────────────────────────
{
\setbeamercolor{background canvas}{bg=redbg}
\begin{frame}[t, plain]
  \begin{tikzpicture}[remember picture, overlay]
    \fill[redacc] (current page.north west)
      rectangle ([yshift=-1.35cm]current page.north east);
    \node[anchor=west, text=white, font=\bfseries\large,
          inner xsep=0.55cm, inner ysep=0pt]
      at ([yshift=-0.68cm]current page.north west)
      {Guided Practice --- Daily Text Messages ($n = 40$)};
  \end{tikzpicture}
  \vspace{1.2cm}\par
  \vspace{0.3em}%
  \hspace{0.55cm}%
  \begin{minipage}{0.9\textwidth}
    \small

    \textit{The dotplot shows: most students send 20--60 messages/day;
    a cluster at 100--140; one student at 180 messages; center near 45;
    range 5 to 180 messages.}

    \vspace{0.6em}
    \begin{enumerate}
      \setlength{\itemsep}{0.5em}
      \item What is the shape of this distribution? Unimodal or bimodal?
            What evidence supports your answer?
      \item Identify any potential outliers. Which values qualify and why?
      \item Estimate the center and describe the spread. Include units.
      \item Write a complete SOCS description in your notes. Label each
            sentence S, O, C, or S in the margin.
    \end{enumerate}
  \end{minipage}
\end{frame}
}

% ─────────────────────────────────────────────────────────────────────────────
% SLIDE 10 — GROUP ACTIVITY
% ─────────────────────────────────────────────────────────────────────────────
{
\setbeamercolor{background canvas}{bg=sky}
\begin{frame}[t, plain]
  \navyheader{Group Activity --- Write \& Critique SOCS Descriptions}
  \vspace{0.8em}%
  \hspace{0.55cm}%
  \begin{minipage}{0.9\textwidth}
    \small

    \textbf{Part 1 (8 min):} Your group receives one distribution (A, B, or C).
    Write a complete SOCS description. Label each sentence.

    \vspace{0.5em}
    \begin{tabular}{@{} >{\bfseries}p{0.12\linewidth} p{0.82\linewidth} @{}}
      A & Ages at a concert ($n = 80$): bimodal peaks at 17--19 and 35--40 yrs;
          no outliers; center $\approx$ 27 yrs; range 14--55 yrs. \\[4pt]
      B & Wait times at a coffee shop ($n = 50$): skewed right; three wait
          $>15$ min (one at 22 min); center $\approx$ 5 min; range 1--22 min. \\[4pt]
      C & Final exam scores ($n = 35$): skewed left; two below 40 pts;
          center $\approx$ 82 pts; range 36--100 pts.
    \end{tabular}

    \vspace{0.6em}
    \textbf{Part 2 (5 min):} Critique this sample response for Distribution B:
    \begin{tcolorbox}[colback=white, colframe=redacc, boxrule=0.6pt, arc=1.5mm,
        left=2mm, right=2mm, top=1mm, bottom=1mm]
    \small\textit{``The distribution is skewed right.
    The center is about 5. The wait times go from 1 to 22.
    There are some outliers.''}
    \end{tcolorbox}
    What is wrong? Rewrite it to earn full AP credit.

    \vspace{0.4em}
    \textbf{Part 3 (3 min):} Individual reflection in your notes.
  \end{minipage}
\end{frame}
}

% ─────────────────────────────────────────────────────────────────────────────
% SLIDE 11 — EXIT TICKET
% ─────────────────────────────────────────────────────────────────────────────
{
\setbeamercolor{background canvas}{bg=sky}
\begin{frame}[t, plain]
  \navyheader{Exit Ticket}
  \vspace{1em}%
  \hspace{0.55cm}%
  \begin{minipage}{0.9\textwidth}
    \small

    A researcher recorded commute times (minutes) for 20 workers.
    The distribution is \textit{skewed right}: most workers commute 10--30 min,
    but a few commute 60--90 min. Center $\approx$ 22 min;
    range 8--88 min; two workers over 75 min stand out.

    \vspace{0.6em}
    \begin{enumerate}
      \setlength{\itemsep}{0.7em}
      \item Write a complete SOCS description of commute times for these
            20 workers. Address all four components with context in every sentence.
      \item A classmate says: ``The distribution is skewed right, so the mean
            commute time is probably \textit{less than} the median.''
            Is the classmate correct? Explain.
    \end{enumerate}

    \vspace{0.5cm}
    {\color{navy}\itshape\small
      Collect exit tickets. Use the quick rubric (1 pt each: S, O, C, S, context)
      to give targeted feedback before Lesson 1.7.}
  \end{minipage}
\end{frame}
}

% ─────────────────────────────────────────────────────────────────────────────
% SLIDE 12 — SPIRAL CONNECTIONS & PREVIEW
% ─────────────────────────────────────────────────────────────────────────────
{
\setbeamercolor{background canvas}{bg=navy}
\begin{frame}[t, plain]
  \begin{tikzpicture}[remember picture, overlay]
    \fill[navylight] (current page.north west)
      rectangle ([yshift=-1.35cm]current page.north east);
    \node[anchor=west, text=white, font=\bfseries\large,
          inner xsep=0.55cm, inner ysep=0pt]
      at ([yshift=-0.68cm]current page.north west)
      {Connections \& Preview};
  \end{tikzpicture}
  \vspace{1.2cm}\par
  \vspace{0.3em}%
  \hspace{0.55cm}%
  \begin{minipage}{0.9\textwidth}

    \begin{columns}[t]
    \begin{column}{0.47\textwidth}
      {\color{goldacc}\bfseries\small Connects Forward}
      \begin{itemize}
        \setlength{\itemsep}{0.45em}\color{white}\small
        \item \textbf{Lesson 1.7:} SOCS center becomes mean or median;
              spread becomes IQR or standard deviation.
        \item \textbf{Lesson 1.9:} Compare two distributions with SOCS ---
              same framework, comparative language.
        \item \textbf{AP Exam:} Describe a distribution is one of the most
              common free-response tasks. SOCS + context earns full credit.
      \end{itemize}
    \end{column}
    \begin{column}{0.47\textwidth}
      {\color{skymid}\bfseries\small Lesson 1.7 Preview}
      \begin{itemize}
        \setlength{\itemsep}{0.45em}\color{white!85!navy}\small
        \item Today: center $\approx$ ``balance point''; spread = rough range.
        \item Next lesson: \textbf{mean}, \textbf{median}, \textbf{range},
              \textbf{IQR}, \textbf{standard deviation} --- exact numbers.
        \item Shape guides which measures to choose. A skewed distribution
              $\to$ median and IQR. Symmetric $\to$ mean and standard deviation.
      \end{itemize}
    \end{column}
    \end{columns}

  \end{minipage}
\end{frame}
}

\end{document}
